\section{Metodología de evaluación}
\subsection{Verificación antes de rendimiento}
Se separan cuatro niveles: referencia matemática, unidad RTL, integración
al pipeline y kernel completo. Cada nivel tiene un oráculo y una puerta de
aceptación. Para la unidad se usarán extremos, ambas mitades, cancelaciones
y vectores aleatorios con semilla; para el pipeline, dependencias y efectos
arquitectónicos únicos. Aprobar una batería finita no equivale a una prueba
formal de todos los estados.

La regresión del core se amplía antes de medir. El software usa un subconjunto
de instrucciones permitido y comprobable. No basta compilar para RV32IC: el
compilador puede emitir instrucciones estándar que el core docente no cubre.
La primera versión de los kernels se escribirá en ensamblador controlado.
Los kernels primarios usarán instrucciones de 32 bits, sin compresión ni
relajación automática, con la misma política en todas las variantes. El
soporte C del core se conservará y verificará por separado; no se introduce
como otro factor de la comparación inicial.

\subsection{Configuraciones predefinidas}
\begin{tabularx}{\textwidth}{@{}lY@{}}
\toprule
\textbf{Factor} & \textbf{Plan inicial, aún no ejecutado}\\\midrule
Kernel & Subtotales INT32 por grupo; después, capa con epílogo completo.\\
Batch & Uno.\\
Dimensiones & \(K\in\{32,128,512\}\), \(N\in\{1,4,16\}\).\\
Grupo principal & \(G=32\).\\
Sensibilidad de grupo & \(G\in\{8,32,128\}\), solo cuando divide \(K\).\\
Semillas & Diez: 20260908 a 20260917; conservar también tensores y hashes.\\
Regímenes de \(z\) & Constante, especializado y leído en ejecución; resultados separados.\\
Memoria inicial & Propuesta de 64 KiB IMEM + 64 KiB DMEM, común a variantes.\\
Microarquitectura & Inicial paralela; dos vías e iterativa en fase posterior.\\\bottomrule
\end{tabularx}\par

\(N=1\) constituye el producto punto; \(N>1\) permite reutilización entre filas.
Casos con colas, por ejemplo \(K\in\{1,3,5,31,33,127\}\), se verifican primero
como funcionalidad. No se incorporan silenciosamente a un agregado de rendimiento
definido para otras configuraciones. La capacidad de memoria propuesta se
confirmará al generar los programas; cualquier cambio será común y registrado.

\subsection{Política de zero-point y baseline factorizado}
\begin{description}[style=nextline]
\item[ZC: constante por ejecución.] Campañas independientes para \(z=0,8,15\).
Todas las variantes conocen el mismo valor antes de ensamblar.
\item[ZS: especializado.] Los valores cambian entre fila/grupo, pero se conocen
antes de generar código. Se cuenta el código adicional y todas las variantes
tienen la misma información para optimizar.
\item[ZR: leído en ejecución.] Los metadatos llegan por memoria. Si D usa
inmediatos, se cuenta su despacho; si cambia el ISA para suministrarlos, se
identifica como variante adicional. Su implementación sigue abierta.
\end{description}

B3 calcula \(S_{a,g}\) una vez por vector y grupo y lo comparte entre filas;
no se fuerza su recálculo para favorecer a D. B3 y D usan el mismo número de
vías y protocolos comparables, aunque sus costos lógicos pueden diferir.
Los dos disponen de MUL escalar para las operaciones restantes.

Los pesos permanecen empacados en U4 en la comparación principal. Una posible
recodificación a S4 cuando \(z=8\) pertenece a otra comparación explícita.
La mejora debida a reducir bytes desde FP32 no se atribuye a una instrucción
cuando el baseline ya recibe los mismos datos U4.

\subsection{Frontera de medición y métricas}
La región comienza con la aceptación de la primera instrucción del kernel y
termina con el retiro de la última que materializa su salida. El testbench
deberá implementar y documentar esos eventos. Se excluyen reset y carga inicial
del programa; se incluyen acceso a tensores, bucles, corrección, suma de
activaciones, despacho de metadatos y stores de salida.

Se medirá aparte el núcleo aritmético y, cuando exista, la capa con epílogo.
El subtotal entero por grupo no se llamará ``inferencia completa''. Las
magnitudes principales son
\begin{equation}
S_{\mathrm{ciclos}}=\frac{C_B}{C_D},\qquad
\mathrm{CPI}=\frac{C}{I_{\mathrm{retiradas}}},\qquad
t=\frac{C}{f}.
\end{equation}
El contador de retiro incluye stores y branches, no solo instrucciones con
escritura de registros. Se conservarán conteos absolutos, stalls desglosados,
código y tráfico por nivel de memoria. Sin timing sustentado, \(S_{\mathrm{ciclos}}\)
no se presentará como speedup físico medido.

\subsection{Análisis y fase física}
Se compararán variantes de forma pareada por configuración, semilla y tensores.
Los casos con \(S<1\) se publican como degradaciones. Cualquier media geométrica
indicará qué cocientes integra; no sustituirá la tabla de casos completos.
La repetición de una simulación determinista con idénticos datos no genera
por sí sola evidencia estadística adicional. Si se calculan intervalos,
se identificará su unidad de muestreo.

Una fase posterior evaluará modelos de memoria con espera y unidades de una,
dos y cuatro vías. Síntesis y place-and-route deberán usar dispositivo,
herramienta y restricciones documentados. Se distinguirá un reloj común
viable de la frecuencia alcanzada por cada diseño. LUT, FF, DSP y BRAM no
son una medida directa de energía; esta requiere potencia y tiempo bajo una
metodología explícita. La FPGA y su acceso universitario siguen pendientes.

\subsection{Reproducibilidad del experimento}
Cada ejecución conservará revisión del core y del protocolo, versiones y
opciones de herramientas, comando, configuración, semilla, hashes del programa
y tensores, salida esperada/obtenida y logs. El esquema inicial vive en
\code{results/schema.json}. No se han creado filas sintéticas de rendimiento.
Un fallo se conserva con su estado; una tabla del artículo se generará desde
los datos, no desde números transcritos a mano.
